\documentclass[11pt,a4paper]{article}
\usepackage[utf8]{inputenc}
\usepackage[T1]{fontenc}
\usepackage{amsmath,amssymb,amsfonts}
\usepackage[margin=2.5cm]{geometry}
\usepackage{hyperref}

\newcommand{\R}{{\mathbb{R}}}
\newcommand{\C}{{\mathbb{C}}}

\title{Diffractive Beam Shaper Design\\via Finite Element Method\\[0.5em]
\large Mathematical Description of the Method}
\author{Pawel Wocjan\\[0.3em]
\small Based on: \emph{Entwurf diffraktiver Strahlformer mit der Methode der finiten Elemente}\\
\small Studienarbeit, IAKS, Universit\"at Karlsruhe (TH), 1997\\
\small Supervised by Prof.\ Dr.\ Th.\ Beth and Dipl.-Inform.\ M.\ Schmid}
\date{}

\begin{document}
\maketitle

\section{The Problem}

Design a \textbf{diffractive phase element} (DPE) that transforms an input beam $f_1$ 
into a desired output intensity pattern~$|f_2|^2$. The DPE is described by its transmission function
\[
  T_{\mathrm{DPE}}(x,y) = \exp\bigl(i\,\phi(x,y)\bigr),
\]
where $\phi(x,y)$ is the phase function to be computed. 
Since $|T_{\mathrm{DPE}}|=1$, no energy is absorbed.

In the optical setup, the input and output planes are at focal distance~$f$ from a 
Fourier-transforming lens. Wave propagation is described by the 2D Fourier transform:
\[
  f_2 = \mathcal{F}\bigl\{T_{\mathrm{DPE}} \cdot f_1\bigr\}.
\]
Since only the \textbf{intensity} $|f_2|^2$ matters (not its phase), 
there is a large design freedom that the algorithms below exploit.

\section{The Method}

\subsection{Step 1: Mesh Generation}

Lay topologically equivalent meshes over the input plane~$E$ and output plane~$A$. 
For \textbf{rectangular topology}, nodes are indexed by $(i,j)$ with $i,j\in\{0,\ldots,s{-}1\}$, 
forming quadrilateral cells. Interior nodes have 4 direct neighbors and 8 surrounding nodes.

Interior node positions are initialized using \textbf{4-point interpolation} (Atkin, 1994):
\[
  P_0 = \sum_{l=1}^{4} g_l\, P_l,
\]
where the weights depend on the Euclidean distances in index space:
\begin{align*}
  g_1 &= \frac{d_2\, d_3\, d_4}{(d_1+d_2)\,(d_1 d_2 + d_3 d_4)}, &
  g_2 &= \frac{d_1\, d_3\, d_4}{(d_1+d_2)\,(d_1 d_2 + d_3 d_4)}, \\[4pt]
  g_3 &= \frac{d_1\, d_2\, d_4}{(d_3+d_4)\,(d_3 d_4 + d_1 d_2)}, &
  g_4 &= \frac{d_1\, d_2\, d_3}{(d_3+d_4)\,(d_3 d_4 + d_1 d_2)},
\end{align*}
with $d_k = \bigl[(i_0-i_k)^2 + (j_0-j_k)^2\bigr]^{1/2}$.

For \textbf{triangular topology}, nodes are indexed by $(i,j)$ with $j=0,\ldots,s{-}1$ 
and $i=0,\ldots,j$, using a 6-point interpolation formula with analogous weights.

\subsection{Step 2: Mesh Optimization}

\subsubsection{Equal-Area Optimization (8-Point Algorithm)}

For constant-intensity signals, equal area implies equal energy. 
For an interior node~$P_0$ surrounded by 8 neighbors $P_1,\ldots,P_8$ 
(numbered counterclockwise from top: $P_1$=top, $P_2$=top-right, \ldots, $P_8$=top-left), 
the \textbf{ideal position} $C(x_c,y_c)$ that equalizes the 4 surrounding cell areas satisfies:
\[
  x_c = \frac{bf - de}{bc - ad}, \qquad y_c = \frac{af - ce}{bc - ad},
\]
where
\begin{align*}
  a &= y_2 - y_8 - y_6 + y_4, & b &= x_2 - x_8 - x_6 + x_4, \\
  c &= y_8 - y_6 - y_4 + y_2, & d &= x_8 - x_6 - x_4 + x_2, \\
  e &= x_1(y_2{-}y_8) + x_5(y_4{-}y_6) + y_1(x_8{-}x_2) + y_5(x_6{-}x_4), \\
  f &= x_3(y_2{-}y_4) + x_7(y_8{-}y_6) + y_3(x_4{-}x_2) + y_7(x_6{-}x_8).
\end{align*}
The node is moved iteratively: $P_0' = P_0 + \alpha\,(C - P_0)$ with step size $0 < \alpha \le 1$.

\subsubsection{Energy-Weighted Optimization (General Signals)}

For signals with non-uniform intensity, the Python implementation uses 
\textbf{energy-weighted Laplace smoothing}. Each interior node is moved toward the 
energy-weighted centroid of its 4 direct neighbors:
\[
  C = \frac{w_{\mathrm{left}}\, P_{\mathrm{left}} + w_{\mathrm{right}}\, P_{\mathrm{right}} 
  + w_{\mathrm{up}}\, P_{\mathrm{up}} + w_{\mathrm{down}}\, P_{\mathrm{down}}}
  {w_{\mathrm{left}} + w_{\mathrm{right}} + w_{\mathrm{up}} + w_{\mathrm{down}}},
\]
where $w_{\mathrm{left}}$ is the total energy in the two cells to the left of the node, etc.
This pulls nodes toward high-energy regions, creating smaller cells there and larger cells 
in low-energy regions. A validity check rejects any move that would create a negative-area 
(tangled) cell.

The energy in each quadrilateral cell is computed by scanning pixels inside the polygon 
using a ray-casting point-in-polygon test.

\subsection{Step 3: Phase Recovery (Polynomial Least Squares)}

The mesh correspondences $\mathrm{in}(i,j)=(x_{ij},y_{ij})$ and 
$\mathrm{out}(i,j)=(\phi_x^{ij},\phi_y^{ij})$ define a discrete geometric transformation. 
By the \textbf{stationary phase condition}:
\[
  \nabla\phi(x_{ij},y_{ij}) = c \cdot 
  \begin{pmatrix} \phi_x^{ij} - \tfrac{1}{2} \\ \phi_y^{ij} - \tfrac{1}{2} \end{pmatrix},
  \qquad c = 2\pi N,
\]
where $N$ is the FFT grid size and the subtraction of $\tfrac{1}{2}$ centers the output coordinates 
(matching the \texttt{fftshift} convention).

We approximate $\phi$ as a bivariate polynomial of degree $D{-}1$:
\[
  \phi(x,y) = \sum_{\substack{k,l \ge 0 \\ k+l \le D-1}} a_{kl}\, x^k y^l,
\]
with $D(D{+}1)/2$ unknown coefficients. The gradient conditions give a linear system:
\[
  M \cdot \mathbf{a} = \mathbf{b},
\]
where the matrix entries are
\[
  M_{\mathrm{row},\mathrm{col}} = \sum_{i,j} \Bigl[
    m\cdot k\cdot x_{ij}^{m-1+k-1}\, y_{ij}^{n+l} +
    n\cdot l\cdot x_{ij}^{m+k}\, y_{ij}^{n-1+l-1}
  \Bigr],
\]
and the right-hand side entries are
\[
  b_{\mathrm{row}} = \sum_{i,j} \Bigl[
    k\cdot \phi_x^{ij}\cdot x_{ij}^{k-1}\, y_{ij}^l +
    l\cdot \phi_y^{ij}\cdot x_{ij}^k\, y_{ij}^{l-1}
  \Bigr],
\]
using the index mapping $\mathrm{row} = l(2D{-}l{+}1)/2 + k$ and 
$\mathrm{col} = n(2D{-}n{+}1)/2 + m$. This system is solved via least squares.

\subsection{Step 4: IFTA Refinement}

The phase from Step~3 serves as a starting point for the 
\textbf{Iterative Fourier Transform Algorithm} (Gerchberg--Saxton/Fienup), 
which alternates between input and output planes:
\begin{itemize}
  \item \textbf{Input plane}: replace amplitude with the known $|f_1|$, keep phase.
  \item \textbf{Output plane}: replace amplitude with desired $|f_2|$ (phase freedom), keep phase.
\end{itemize}
The Studienarbeit runs 10 iterations for efficiency, then 20 for SNR.

\subsection{The SepOp Baseline}

The Separabilisierungsoperator maps any 2D signal to a nearby separable signal:
\[
  \mathrm{SepOp}\colon f(x,y) \;\mapsto\; 
  \underbrace{\left(\int f(\xi,y)\,d\xi\right)}_{p(y)} \cdot 
  \underbrace{\left(\int f(x,\xi)\,d\xi\right)}_{q(x)}.
\]
In the discrete case: $q(x) = \sum_y f(x,y)$ and $p(y) = \sum_x f(x,y)$, 
giving $s(x,y) = q(x)\cdot p(y)$. For separable signals, the 2D problem decomposes 
into two independent 1D problems with analytical solutions.

\section{Quality Metrics}

\textbf{Signal-to-Noise Ratio} (SNR):
\[
  \mathrm{SNR}_I(f_2,f) = \frac{\|f\|^2}{\left\| |f_2| - \alpha|f| \right\|^2},
  \qquad \alpha = \frac{\langle |f|,\, |f_2|\rangle}{\|f\|^2}.
\]

\textbf{Diffraction Efficiency}:
\[
  \eta = \frac{\|f\|^2_{W_{\mathrm{Signal}}}}{\|f_1\|^2_{W_{\mathrm{DE}}}}.
\]

\section{Original Results (1997)}

\begin{center}
\begin{tabular}{|l|l|r|r|}
\hline
\textbf{Input $\to$ Output} & \textbf{Method} & \textbf{SNR (dB)} & $\boldsymbol{\eta}$ \textbf{(\%)} \\ \hline\hline
Circle $\to$ Marilyn & SepOp  & 19.3 & 78.6 \\
                     & Kugel  & 16.9 & 75.1 \\
                     & Random & 13.8 & 68.9 \\
                     & \textbf{FEM} & \textbf{20.8} & \textbf{79.0} \\ \hline
Insep.\ $\to$ Marilyn & SepOp  & 32.8 & 92.1 \\
                       & \textbf{FEM} & \textbf{35.8} & 86.0 \\ \hline
Gauss $\to$ Marilyn & \textbf{SepOp} & \textbf{31.4} & \textbf{92.5} \\
                    & FEM   & 29.9 & 83.0 \\ \hline
Quadrat $\to$ Marilyn & SepOp & 20.8 & 79.1 \\
                      & FEM   & 20.5 & 78.0 \\ \hline
\end{tabular}
\end{center}

The FEM method excels for non-separable input signals and performs comparably 
to SepOp for separable ones.

\section{References}

\begin{enumerate}
\item H.~Aagedal, M.~Schmid, S.~Egner, J.~M\"uller-Quade, Th.~Beth, F.~Wyrowski,
  \emph{Analytical beam shaping with application to laser diode arrays}, 
  JOSA~A, Vol.~14, No.~7, 1997.
\item T.~Dresel, M.~Beyerlein, J.~Schwider, 
  \emph{Design of computer-generated beam-shaping holograms by iterative 
  finite-element mesh adaptation}, Applied Optics~35, 6865--6874, 1996.
\item R.W.~Gerchberg, W.O.~Saxton, 
  \emph{A practical algorithm for the determination of phase from image and 
  diffraction plane pictures}, Optik~35, 237--264, 1972.
\item F.~Wyrowski, O.~Bryngdahl, 
  \emph{Iterative Fourier-transform algorithm applied to computer holography}, 
  JOSA~A, Vol.~5, No.~7, 1988.
\end{enumerate}

\end{document}
